\subsection{LLM-Assisted Background Knowledge Elicitation in the Custom ILP Pipeline}
\label{sec:custom-ilp-llm}

Beyond the fixed benchmark datasets, the platform includes a
\textit{Custom ILP} tab that lets a user upload an arbitrary tabular
dataset (e.g.\ a spreadsheet of patients and symptoms) and convert it,
automatically, into a genuine Popper or Andante learning problem. Each
row becomes a sample identifier, each feature column becomes one or
more unary (or binary) background predicates, and the user-selected
target column and positive class become the positive/negative
examples. This conversion pipeline is described separately
(Section~\ref{sec:tabular-to-ilp}); this section documents the large
language model (LLM) layer added on top of it, whose purpose is to let
the user express \emph{domain knowledge} in free text and have it
compiled into background-knowledge clauses, and to let the user ask
natural-language questions about the resulting dataset and hypothesis.

\paragraph{Motivation.} A tabular-to-predicates conversion captures only
what is literally present in the columns of the uploaded table. A
clinician uploading a patient/symptom table, however, typically also
holds implicit domain rules that are never encoded as a column --
``fever and cough together suggest flu'', for instance. Without a way
to inject such rules, the learner is restricted to combining raw
column predicates directly, at the cost of both search space and of
discoverable hypotheses that a domain expert would consider obvious.
The platform therefore exposes a free-text box in which the user can
write such rules in either English or French, in an ``if \ldots\ then
\ldots'' / ``si \ldots\ alors \ldots'' style, which the system compiles
into additional Prolog background clauses appended to \texttt{bk.pl}
(Popper) or the \texttt{begin\_bg}/\texttt{end\_bg} block (Andante).

\paragraph{Two-stage design: keyword parsing before LLM parsing.} The
feature was deliberately developed in two stages rather than starting
directly from an LLM-based parser. The first stage
(\texttt{core.tabular\_to\_ilp.parse\_background\_text}) is a
regex-based parser matching a small family of if/then sentence
patterns in both languages, followed by fuzzy substring matching of
each clause of the sentence against the set of predicate names already
generated by the conversion pipeline for that dataset. This stage has
no external dependency, no latency, no cost, and is fully
deterministic -- and it was kept as the default, always-available
fallback after the LLM stage was added, rather than being replaced by
it. The LLM stage
(\texttt{core.llm\_background.parse\_background\_text\_with\_llm}) is
an opt-in enhancement, enabled by a checkbox in the interface, intended
to handle the far larger space of phrasings a keyword parser cannot
cover (synonyms, indirect phrasing, multi-clause conditions), at the
cost of introducing a third-party API dependency and the general
reliability concerns that come with generative text output.

\paragraph{Provider choice.} The LLM integration was first implemented
against the Anthropic Messages API, then migrated to Google's Gemini
API (\texttt{google-generativeai} client library, model
\texttt{gemini-3.8-flash}) once a free usage tier requiring no billing
setup was identified as a hard requirement for the platform: unlike
the Anthropic API, a Gemini API key can be obtained at no cost from
Google AI Studio and used directly. Only the client-facing module
(\texttt{core/llm\_background.py}) needed to change; the pipeline's
safety contract, described next, was preserved unmodified across the
migration, since it does not depend on which provider produced the
text.

\paragraph{Safety contract: the model never writes Prolog directly.}
The central design constraint of this feature is that the LLM's raw
output is never trusted as executable logic. Concretely:
\begin{enumerate}
\item The model is prompted with the exact list of predicate names
already known for the current dataset (i.e.\ those actually generated
by the tabular-to-ILP conversion, Section~\ref{sec:tabular-to-ilp}) and
instructed to answer using \emph{only} names drawn verbatim from that
list, or \texttt{null} where no known predicate corresponds to a
concept in the user's text -- never to invent a new predicate name.
\item The model is constrained to return a single JSON array (via the
API's structured-output mode, \texttt{response\_mime\_type:
"application/json"}), each element of the form
\texttt{\{"conclusion": \ldots, "conditions": [\ldots]\}}, with no
surrounding prose.
\item After parsing this JSON, \emph{every} conclusion and every
condition string returned by the model is independently re-checked
against the same known-predicate set before any Prolog text is
constructed. A rule referencing even one unrecognized predicate name
is discarded in full and reported to the user as a warning, rather
than partially compiled.
\item Only once a rule's conclusion and all of its conditions have
passed this check is the corresponding clause emitted, in the fixed
template \texttt{conclusion(S) :- cond\_1(S), \ldots, cond\_n(S).},
where \texttt{S} is the shared sample-identifier variable.
\end{enumerate}
This validate-after-generate pattern -- identical in spirit to
treating any language-model output as untrusted user input --
guarantees that a hallucinated or slightly misspelled predicate name
can never reach the generated \texttt{bk.pl}, regardless of what the
model returns; at worst, a legitimate rule the user intended is simply
dropped and surfaced as a warning rather than silently miscompiled.
Any failure of the call itself (missing or invalid API key, network
error, malformed or non-JSON response) is caught and raised as a
dedicated \texttt{LLMParsingUnavailable} exception, which the calling
page catches to fall back transparently to the keyword-based parser
of the first stage, so that a transient API failure degrades the
feature rather than breaking the pipeline.

\paragraph{Conversational interface over the dataset.} A second,
independent use of the same client library is a chat interface
(\texttt{core.llm\_background.chat\_reply}) allowing the user to ask
free-form questions about the uploaded dataset, the generated
predicates, and -- once a run has completed -- the learned hypothesis
and its accuracy. Unlike the background-knowledge parser, this
component's output is never compiled into logic and carries no
correctness requirement beyond being a helpful natural-language
answer; its context window is constructed from the conversion plan
(feature predicates and their human-readable glossary), the dataset
summary, and, when available, the stored result of the most recent
run, all assembled server-side and passed as a system instruction
before the model sees the user's question. Two variants of this
chat were implemented: one scoped to a dataset already uploaded
(``discuss this specific data''), and a second, upload-independent
variant offered before any file is selected, intended to help a user
who is unsure what shape of table or target column the tool expects.

\paragraph{Operational notes.} The Gemini API key is read at runtime
from \texttt{.streamlit/secrets.toml}, a file explicitly excluded from
version control; both the background-knowledge checkbox and the chat
boxes detect its absence and disable themselves accordingly rather
than raising an error, so the platform remains fully usable --
including its core ILP functionality -- without any LLM credentials
configured. No user data or generated predicate is transmitted to the
API beyond what is needed to answer the specific request (the known
predicate list, the user's free text, and, for the chat, the
in-session conversation history).
